% \section{Discussion and Outlook}
% \label{Discussion:Outlook}

% We have studied the Kitaev-Heisenberg model on the star lattice and determined the ground state phase diagram. By tuning the relative strength of the Kitaev and Heisenberg couplings, we unearth a rich phase diagram with six distinct phases. Two of these -- the FM and stripy order -- have analogues on the corresponding honeycomb  model. Extended chiral QSL phases are found around the exact Kitaev limits, with the FM QSL being more stable than the AFM one. In both phases, we find non-zero Chern numbers of the effective Majorana bands which originate from the spontaneous breaking of TRS due to fluxes on the triangular plaquettes  of the star lattice. We confirmed the existence of chiral Majorana edge modes weakly renormalized by the presence of the Heisenberg interaction.

% In the AFM range of the phase diagram we find two phases with no analogues on the honeycomb lattice, i.e. a singlet VBS and an exotic triplet VBS. These appear because the simple Neel AFM ordering that appears on the Honeycomb is geometrically frustrated on the star lattice. We unveil a symmetry between the two VBS states by generalizing the sublattice-dependent spin rotation which relates the Hamiltonian of different parts of the phase diagram. 
% We then develop a bond operator formalism to show that the VBS phases have  a rich topological structure for the triplon excitations. The  introduction of a magnetic field in conjunction with the Kitaev interaction leads to gapped bands with a whole variety of (large) Chern numbers. 
% Accordingly we find that the Kitaev-Heisenberg model on the star lattice cannot only support chiral Majorana excitations but  also chiral triplon edge modes.  


% Our work raises a number of questions for further study. First, the star lattice provides one of the simplest examples of a non-bipartite lattice. For the pure Kitaev model the presence of odd-number triangle plaquettes led to the first exactly soluble \textit{chiral} QSL with spontaneously broken  TRS~\cite{yao_exact_2007}, and for the pure Heisenberg limit was shown to realize a VBS ground state~\cite{richter_absence_2004}. Thus, an interesting avenue will be to explore whether the Kitaev-Heisenberg model shows similarly rich phase diagrams on other non-bipartite decorated lattices, possibly with exotic VBS phases. Second, the new triplet VBS phase we discover around $\varphi \approx 1.7\pi$ shows intriguing topological triplet and singlet excitations but many open questions remain. For example, in the absence of an external magnetic field the dynamics are expected to be identical to the singlet VBS phase due to the four-unit-cell rotation discussed above. However an applied magnetic field breaks this symmetry, and the possibility remains that this phase could have qualitatively very different properties than the singlet VBS. In addition, it will be important to understand the stability of the topological triplons and their experimental signatures.
% Third, we showed that a sublattice dependent spin rotation can connect different parameter regimes of the Kitaev-Heisenberg model similar to the honeycomb lattice. Here on the star lattice it also uncovers a connection between two distinct VBS phases. It will be interesting to explore other lattice generalizations of such isometries, which would possibly allow us to find exact magnetic ground states even on amorphous lattices with net zero magnetization. 

% Finally, there is a rich variety of modifications to the Hamiltonian that could be made in order to further assess the robustness of the different phases and, especially their chiral excitations, as described here. This is particularly timely in light of recent experimental work which demonstrated the successful synthesis of a spin $1/2$ star lattice structure with antiferromagnetic couplings~\cite{sorolla_synthesis_2020}. The material is expected to have in-triangle couplings much stronger than the inter-triangle couplings. Thus it would be worth considering the phase diagram when the relative strength of couplings is adjusted. In addition, it would be interesting to derive microscopic Hamiltonians based on the microscopic orbital configuration and ab-initio calculations, e.g.~considering the effect of including symmetric off-diagonal exchanges~\cite{rau_generic_2014, knolle_dynamics_2018}, as well as Dzyaloshinskii-Moriya interactions~\cite{kim_realization_2016, zhang_interplay_2021}.

% In conclusion, the two limits of the pure Kitaev and Heisenberg models on the star lattice have been known to realize sought-after quantum magnentic phases, e.g.~a chiral QSL and a singlet VBS. We have shown that the phase diagram of the full Kitaev-Heisenberg model is much richer and allows to tune from a phase with chiral Majorana excitations to an exotic bond-selective triplet VBS phases with chiral triplon excitations. We expect that  competing exchange interactions on other decorated lattices can lead to similarly rich physics and hope that material synthesis will help with their realisation. 